\honest{Toolbox-thinking and Law-thinking}{Eliezer Yudkowsky}{2018}

There are two styles of thought. Toolbox thinkers keep a big bag of methods and choose among them by context, and suspect that anyone who talks of a single optimal way does not know the other tools. Law thinkers, done right, separate descriptive truths, normative ideals and prescriptive advice; they may call a path optimal even if no practical algorithm finds it, and they recognize truths that are not tools. In a nearly flat maze, the triangle inequality, that path AC is never longer than path ABC, is always true but only sometimes useful. It tells you to take AC only if you know which path is which and care only about distance, not, say, about avoiding stairs. But it governs the maze either way. To a toolbox thinker, badly explained, this sounds like ``Throw away all the other tools in your toolbox!'' \nb{The post introduces the toolbox style by its suspicion of other people, and the law style as it works ``done correctly.''}

My occasion is an exchange on Twitter. David Chapman defined rationalism as any claim of an ultimate criterion for correct thinking and acting, going beyond ``systematic methods are often useful, hooray!'': a belief in ``one weird trick'' to correct thinking, and added that some rationalisms ``insist there must be one, but that it is not currently knowable.'' Julia Galef replied that rationalists hold one correct normative model, which no single rule can approximate in practice, and look for collections of tricks that bring them closer, such as taking more small risks to raise expected value. Chapman said that understanding rationality as a bag of tricks is what is ``meta-rational.'' Galef then quoted a paper that separates the ideal from practical advice. Chapman answered that this distinction, which he credited to Jonathan Baron, was new to him, ``useful and maybe key,'' and that the real disagreement might be whether an ideal ``that can't be achieved, even in principle, is a good thing.''

``I'm now going to badly stereotype this conversation in the form I feel like I've seen it many times previously.'' And I do. Msr.\ Toolbox says that the many ways of computing p-values are a family of tools, each good in some contexts; likelihood ratios surely have uses, but it would be surprising if one trick were best in every paper, and claiming so shows ``the sheer folly of callow youth'' with only a hammer. Msr.\ Lawful says that even when we cannot compute exact Bayesian updates, the math describes the optimal update, as a Carnot cycle describes an ideal engine no one can build, and that many coherence theorems say that behavior incoherent with probability theory is a dominated strategy. My stereotyped toolbox thinker hears this as a recipe to execute everywhere, and cannot see what else an ``ideal'' algorithm could be for. \nb{Chapman, in his last reply, did not; he asked whether an ideal that cannot be reached is worth having. The post answers that only through a maze.} I allow that there may be ``a symmetrical form of durable misunderstanding'': a law thinker, told that a recipe is just a useful tool, may not see what a recommended recipe could be except a claim to be right, and may suspect ``some absurd motte-and-bailey.'' And the callow youths with a One True Recipe, but no ideal to prove it optimal, do exist. The confusion resolves by grasping a distinction, as a lawful person would put it, or by using both kinds of thinking by context, as a toolbox person would.

Then comes a long maze. Every maze has a shortest path. The fact that a straight path is never longer than a detour always holds, but it tells you which way to walk only if you know the paths and care only about distance, and not, say, about avoiding stairs because someone in your party uses a wheelchair. Its existence does not let you pick the shorter branch at each turn, and that does not make distance useless. A walker with a poor memory may do better keeping the left hand on the wall, even when told a right turn stays on the shortest path, since jumping walls just once can strand you on a disconnected island of walls. That rule gives the shortest expected walk given your resources, yet you may not be on the shortest path. Knowing that a shortest path exists shows that there is room to improve, perhaps with a maze-mapping phone app. The room exists not because no recipe is perfect: given a map with the shortest path drawn, you could walk it. Shortness is a property of paths; a tendency to produce short paths is a property of recipes. The app is better because its path is shorter in a way defined apart from any algorithm. Once you admit ``shorter,'' it is hard not to admit ``shortest,'' and refusing to think about ideal solutions would remove a useful tool from your toolbox. Why would you? \nb{But the maze comes with an agreed measure, distance, and a map would show the shortest path. Whether the ideal for thinking can be reached even in principle is what Chapman asked. The maze assumes that it can.}

Next, the objections to law thinking. I write them myself, saying I am not sure I could state the toolbox side's view to its satisfaction, and I write them as sweeping claims: Msr.\ Law ``will inevitably forget'' the context, and it ``will also never occur'' to Msr.\ Law to ride a bicycle round the maze. The objections also say that people will confuse their current recipe with the unseen ideal and use the ideal to praise it, that Msr.\ Law will send a wheelchair user down the ``best'' path with its stairs, and that a helicopter over the maze will break Msr.\ Law's ontology. Among them is the one real case in the post: law thinkers who ``assume that the behavior of mortgage-backed securities is independently Gaussian-random because the math is neater,'' and bring down their trading firm. My reply is, ``Well, yes, that happens some of the time,'' and sometimes Msr.\ Law needs a lecture on diverse toolboxes; but not everyone needs it, and extreme toolbox thinkers could use a lecture on unreachable ideals and the obstacles between us and them. \nb{The post wrote the objections as if these failures always happen, so ``some of the time'' defeats them. It never discusses the mortgages. The model usually blamed, David Li's Gaussian copula, described how defaults depend on each other; it did not assume they were independent.}

My resolution is that both styles are ``metatools in the metatoolbox,'' to be used ``in ways that depend on context and circumstance.'' Earlier I had given this answer in two versions, a lawful one and this one, ``as a more Toolbox-inclined person might put it.'' Thinking in laws is often useful if you understand the context and caveats. That is not the same as saying every law has exceptions: thermodynamics still holds while you play tennis, when it is no time to think about it. A healthy law thinker will suspect there are laws for how to contextualize, even unknown ones, and meanwhile use chaotic-seeming prescriptions without mistaking them for laws or taking their chaos as proof that no ideal exists.

I close with a remedy for toolbox thinkers: see laws as descriptive statements, not ideals. Once distance is defined, a shortest path simply exists. The rule ``try walking along AC'' holds only if you want shorter walks, know which turn is which, are avoiding no stairs and own no bicycle; the triangle inequality holds regardless. That no heat engine beats a Carnot cycle is not a hymn to the Carnot cycle but a fact, which may suggest looking at heat loss to improve an engine, and which no engine escapes when you are not thinking about it. The best examples of laws are plain facts: that a straight path is never longer than a detour, the limit on the efficiency of engines, and Bayesian updating, since no different update can get a higher expected score. \nb{That last holds only if the prior probabilities and likelihoods are right, and the mortgage models are an example of when they were not.} You cannot get more evidence from an observation than its likelihood ratio gives, and seeing this as a law over all recipes, Bayesian or not, is ``a deeper level of understanding.'' Likewise the deeper reading of the coherence theorems is that incoherence of kind X corresponds to a dominated strategy of kind Y, true of every statistical tool, whether or not you care or can compute better, ``even if you own a bicycle.''

In a clause near the end I concede that for values, such as the ideal of a good universe, ``there isn't any deeper fact of nature.'' \nb{Chapman's definition was about correct thinking and acting. The laws for acting that the post cites judge an action only against the actor's own values.}
